\section{Complete equality classification}\label{sec:equality}
\begin{theorem}[All joint-frontier maximizers]\label{thm:joint-equality}
For every $0\le c\le1$ and $0\le b<1$, a two-qubit state attains \eqref{eq:joint} if and only if it is locally unitarily equivalent to \eqref{eq:joint-state}.
\end{theorem}
\paragraph{Idea of the proof.} Equality must hold in every step of the
proof of Theorem~\ref{thm:weighted}. The extracted vector $w$ must then have
the largest admissible Schmidt coefficients, the positive residual
$\sigma-\ket w\bra w$ must live where the target cannot see it, and the filter
must be aligned with the Schmidt bases of $w$. Each of these conditions is
rigid, and together they determine the state up to local unitaries.

\begin{proof}
First assume $0<c<1$. Choose a maximizing maximally entangled $e$ and write the state as a filter of its canonical $\sigma$. Put $x=(\Id\otimes K^\dagger)e$, so $x$ is unit with Schmidt amplitudes
\[
 s_1=\sqrt{(1+b)/2},\qquad s_2=\sqrt{(1-b)/2}>0.
\]
Extract $w=\sigma x/\sqrt{t}$, $t=\bra x\sigma\ket x=\Fe(c,b)>0$. In the proof of Theorem~\ref{thm:weighted} every inequality must be an equality. Both positive $s_i$ force the squared Schmidt amplitudes of $w$ to be
\[
 q_1=\frac12,\qquad q_2=\frac{1-c}{2}.
\]
Apply local unitaries to put
\[
 w=\frac1{\sqrt2}\bigl(\ket{00}+\sqrt{1-c}\ket{11}\bigr),\qquad
 R=\sigma-\ket w\bra w\succeq0.
\]
The Bob marginal then satisfies
\[
 R_B=\frac c2\ket1\bra1.
\]
A positive operator with its Bob marginal supported on $\ket1$ is supported on $\C^2\otimes\mathrm{span}\{\ket1\}$: the expectation on every vector with Bob component $\ket0$ is zero, and a positive operator annihilates a vector of zero expectation. Hence $R=H\otimes\ket1\bra1$ with $H\succeq0$ and $\tr H=c/2$.

Now $\sigma_A=\diag(q_1,q_2)+H$ has smallest eigenvalue $q_2$. Take a unit eigenvector $u$ for that eigenvalue. The two nonnegative terms
\[
 \ip{u}{(\diag(q_1,q_2)-q_2\Id)u}\quad\text{and}\quad\ip{u}{Hu}
\]
sum to zero. Because $q_1>q_2$, $u$ is a phase multiple of $\ket1$, and positivity implies $H\ket1=0$. Therefore $H=(c/2)\ket0\bra0$. We have proved that the canonical state is exactly
\begin{equation}\label{eq:sigma-ad}
 \sigma_c=\frac12\bigl(\ket{00}+\sqrt{1-c}\ket{11}\bigr)
                    \bigl(\bra{00}+\sqrt{1-c}\bra{11}\bigr)
            +\frac c2\ket{01}\bra{01}.
\end{equation}

It remains to check the filter orientation, rather than merely its singular values. In the chosen Schmidt bases write $W=\diag(w_1,w_2)$ with $w_1>w_2>0$, and let $X$ be the coefficient matrix of $x$. Then
\[
 |\tr(W^\dagger X)|
 \le(w_1-w_2)|X_{00}|+w_2|\tr X|
 \le(w_1-w_2)s_1+w_2(s_1+s_2).
\]
Equality forces $|X_{00}|=s_1$. Since $s_1$ is the operator norm of $X$, the off-diagonal entries in its first row and column must vanish. The remaining diagonal entry has modulus $s_2$, and equality in the triangle inequality aligns the two phases. Thus, up to an irrelevant overall phase,
\[
 x=s_1\ket{00}+s_2\ket{11}.
\]
This argument also covers $b=0$, when $s_1=s_2$; only the strict inequality $w_1>w_2$ was used.

If $E$ is the coefficient matrix of the maximally entangled $e$, then $E^\dagger E=\Id/2$ and $X=E\overline K$. Consequently
\[
 K^\dagger K=2(X^\dagger X)^T=\diag(1+b,1-b).
\]
Polar decomposition absorbs the remaining output unitary into a local Bob unitary. Filtering \eqref{eq:sigma-ad} by this positive diagonal matrix's square root gives exactly \eqref{eq:joint-state}.

For $c=0$, equality in both positive-weight caps gives $\|w\|^2=1$. The positive residual has trace zero and vanishes, so $\sigma$ is a pure maximally entangled state. The filtered state is pure, and its Schmidt spectrum is fixed by $b$; all such pure states form one local-unitary orbit, represented by $\rho_{0,b}$.

For $c=1$, $\sigma_A$ is pure, so $\sigma=\ket0\bra0\otimes\Id/2$ up to local unitaries. The state $\rho$ is pure Alice tensor Bob, and all states with the given $b$ again form one local-unitary orbit. Their FEF is $(1+b)/4=\Fe(1,b)$. The converse in all cases follows from the explicitly checked attaining family.\end{proof}

\begin{corollary}[All saturators of Milne's original inequality]\label{cor:original-equality}
For each $c\in[0,1]$ there is exactly one local-unitary orbit of states satisfying $f(\rho)=1-c/2$. A representative is
\begin{equation}\label{eq:milne-saturator}
 \boxed{\rho_c=
 \frac{\ket{\chi_c}\bra{\chi_c}+c(1-c)\ket{01}\bra{01}}{2-c},
 \qquad\ket{\chi_c}=\ket{00}+(1-c)\ket{11}.}
\end{equation}
At $c=0$ these are exactly pure maximally entangled states; for $0<c<1$ they have rank two; at $c=1$ they are exactly pure product states. No full-rank state saturates.
\end{corollary}
\begin{proof}
For $b<1$, combine \eqref{eq:deficit} with Theorem~\ref{thm:joint-equality}: equality forces $b=c/(2-c)$ and substitution gives \eqref{eq:milne-saturator}. This includes $c=0$ and excludes $c=1$ within the full-rank Bob stratum.

For $b=1$, $\rho=\rho_A\otimes\ket\beta\bra\beta$. A direct Schmidt/Bell calculation gives
\begin{equation}\label{eq:singular}
 f(\rho)=\frac12\lambda_{\max}(\rho_A)=\frac{1+c}{4}.
\end{equation}
Equality with $1-c/2$ forces $c=1$, hence both marginals are pure.\end{proof}

\subsection{What happens at a singular Bob marginal}
For fixed $c<1$, $\lim_{b\to1^-}\Fe(c,b)=1/2$, whereas the exact value at $b=1$ is $(1+c)/4<1/2$. There is no discontinuity of the physical function $f(\rho)$: the matrices $\rho_{c,b}$ tend to the pure product $\ket{00}\bra{00}$, whose actual singular centre is $c=1$. It is the coordinate $c$ that fails to extend continuously along this approach. Replacing the singular branch by the limit would answer a different question.

